%--------------------------------------------------------------------%
% PETA KEWAJIBAN PER TAHAP
% Halaman ini otomatis hilang bila \petunjukfalse di thesis.tex.
%--------------------------------------------------------------------%
\chapter*{PETA KEWAJIBAN PER TAHAP}
\thispagestyle{capstone}
{\itshape\textbf{Seminar proposal:} dilaksanakan sebelum pelaksanaan; yang dinilai adalah kesiapan dan kematangan rencana. \textbf{Sidang akhir:} dilaksanakan setelah luaran dievaluasi mitra; yang dinilai adalah capaian dan keterpenuhan tujuan proyek.\par}

{\footnotesize\singlespacing
\wajib{WAJIB}: harus ada. \wajib{WAJIB bila berlaku}: bila kondisinya terpenuhi. \opsional{Opsional}: boleh dikosongkan. \tidakperlu{Tidak diperlukan}: tidak diisi.\par}

\begingroup
\footnotesize\singlespacing
\renewcommand{\arraystretch}{1.35}
\begin{longtable}{|>{\raggedright\arraybackslash}p{0.44\textwidth}|>{\centering\arraybackslash}p{0.22\textwidth}|>{\centering\arraybackslash}p{0.22\textwidth}|}
	\hline
	\rowcolor{blue!12}\textbf{Bagian} & \textbf{Seminar proposal} & \textbf{Sidang akhir}\\ \hline
	\endfirsthead
	\hline
	\rowcolor{blue!12}\textbf{Bagian} & \textbf{Seminar proposal} & \textbf{Sidang akhir}\\ \hline
	\endhead
	P.1 Profil dan legalitas mitra & \cW{WAJIB} & \cT{Tidak diperlukan}\\ \hline
	P.2 Tim dan pembagian kerja & \cW{WAJIB} & \cT{Tidak diperlukan}\\ \hline
	P.3 Rencana pelaksanaan & \cW{WAJIB} & \cT{Tidak diperlukan}\\ \hline
	Pernyataan penggunaan AI (tiap anggota) & \cT{Tidak diperlukan} & \cW{WAJIB}\\ \hline
	Kata pengantar & \cO{Opsional} & \cW{WAJIB}\\ \hline
	Abstract dan Abstrak & \cW{WAJIB (draf)} & \cW{WAJIB (final)}\\ \hline
	BAB I Masalah mitra dan tujuan proyek & \cW{WAJIB (draf)} & \cW{WAJIB (final)}\\ \hline
	BAB II Analisis dan perancangan & \cW{WAJIB (draf)} & \cW{WAJIB (final)}\\ \hline
	BAB III Pelaksanaan proyek & \cT{Tidak diperlukan} & \cW{WAJIB}\\ \hline
	BAB IV Implementasi prototipe & \cO{Opsional (prototipe awal)} & \cW{WAJIB (final)}\\ \hline
	BAB V Pengujian dan evaluasi & \cT{Tidak diperlukan} & \cW{WAJIB}\\ \hline
	BAB VI Penerimaan mitra & \cT{Tidak diperlukan} & \cW{WAJIB}\\ \hline
	BAB VII Kontribusi individu dan refleksi & \cT{Tidak diperlukan} & \cW{WAJIB}\\ \hline
	BAB VIII Kesimpulan dan saran & \cT{Tidak diperlukan} & \cW{WAJIB}\\ \hline
	Kerahasiaan data mitra & \cT{Tidak diperlukan} & \cW{WAJIB}\\ \hline
	Daftar pustaka & \cW{WAJIB bila berlaku} & \cW{WAJIB bila berlaku}\\ \hline
	Lampiran 1 sampai 3 (legalitas, kesediaan, pembagian kerja) & \cW{WAJIB} & \cW{WAJIB}\\ \hline
	Lampiran 4 dan 5 (evaluasi mitra, Form Selesai TA) & \cT{Tidak diperlukan} & \cW{WAJIB (Lampiran 5 bila mitra tidak hadir)}\\ \hline
	Lampiran 6 sampai 9 (\textit{logbook}, \textit{peer assessment}, TKT, kode dan peta kode) & \cO{Opsional (logbook, kode)} & \cW{WAJIB}\\ \hline
	Lampiran 10 (linimasa) & \cW{WAJIB (kolom rencana)} & \cW{WAJIB (lengkap)}\\ \hline
	Lampiran 11 (berita acara seminar proposal) & \cW{WAJIB} & \cT{Tidak diperlukan}\\ \hline
	Lampiran 12 (jalur kompetisi) & \cT{Tidak diperlukan} & \cW{WAJIB bila berlaku}\\ \hline
	Lampiran 13 (pembagian bagian individu dan giliran presentasi) & \cO{Opsional} & \cO{Opsional (disarankan)}\\ \hline
	Lampiran 14 sampai 16 (Form Kendali Bimbingan, \textit{similarity}, log AI dan riwayat versi) & \cW{WAJIB (Form Kendali Bimbingan)} & \cW{WAJIB}\\ \hline
\end{longtable}
\endgroup
{\footnotesize\itshape Pembagian ini disusun dari Pedoman TA Prodi IF dan rubrik TA non-konvensional; konfirmasikan ke Koordinator TA bila ragu.\par}
\clearpage
